\begin{python0}
from solutions import *; clear()
\end{python0}
\sectionthree{``Why do I sometimes see Parameter Names for Function Prototypes?''}

I already told you that function prototypes only need the name of the
functions and the various types: the parameter types and return types.

However sometimes you see parameter names. For instance instead of this
prototype:

\begin{consolethree}[escapeinside=||]
int getSalary(int, int);
\end{consolethree}

you might have this:

\begin{consolethree}[escapeinside=||]
int getSalary(int employeeId, int overtime);
\end{consolethree}

The reason for giving parametric names in function prototypes even
though they are ignored by C++ is that they make the function prototype
\EMPHASIZE{easier to read and use}. This is especially the case where the
prototype has many parameters of the same type. For instance this

\begin{consolethree}[escapeinside=||]
int getSalary(int employeeId, int overtime);
\end{consolethree}

is more useful than this:

\begin{consolethree}[escapeinside=||]
int getSalary(int, int);
\end{consolethree}

Giving parameter names also allows the programmer to document the
function:

\begin{consolethree}[escapeinside=||]
//---------------------------------------------------
// The getSalary() function accepts employeeId and
// overtime (in number of minutes) and returns the
// salary (in cents). Note that a valid employeeId
// ranges from 10000 to 99999.
// If the employeeId is not valid or if the overtime
// is negative the return value is -1.
//---------------------------------------------------
int getSalary(int employeeId, int overtime);
\end{consolethree}

